% Source: PDF pp. 2, 4-5. The downloaded watermark is not book content.
\chapter*{Nguồn, bản quyền và ghi công}
\begingroup\small
\emph{Category Theory for the Sciences}, David I. Spivak, The MIT Press,
Cambridge, Massachusetts và London, Anh, 2014.

\noindent Bản quyền bản gốc: © 2014 Massachusetts Institute of Technology.
ISBN bản bìa cứng: 978-0-262-02813-4.

Ấn bản PDF được dùng làm nguồn vẫn giữ thông báo bản quyền của lần xuất bản
ban đầu: mọi quyền được bảo lưu; không phần nào của sách được sao chép dưới bất
kỳ hình thức nào bằng phương tiện điện tử hoặc cơ học, kể cả sao chụp, ghi âm
hay lưu trữ và truy xuất thông tin, nếu chưa có sự cho phép bằng văn bản của nhà
xuất bản.

\paragraph{Điều kiện sử dụng của ấn bản truy cập mở.}
Khi phát hành sách theo hình thức truy cập mở, MIT Press công bố giấy
phép Creative Commons Ghi công - Phi thương mại - Chia sẻ tương tự 4.0 Quốc tế
(CC BY-NC-SA 4.0) cho ấn bản này:
\url{https://mitpress.mit.edu/9780262028134/category-theory-for-the-sciences/}.
Thông báo cấp phép đó cũng có trong trang bản quyền của bản HTML do nhà xuất
bản phát hành, được lưu trong hồ sơ nguồn của dự án.

\paragraph{Thông tin về bản dịch.}
Bản dịch máy tiếng Việt này sử dụng AI \textbf{GPT-6.1 Sol (Ultra)}.
Người thực hiện prompting: \textbf{Bế Trọng Nghĩa} (\nolinkurl{@ichxorya}).
Bản dịch được dàn trang lại và Việt hóa nhãn sơ đồ. Đây là một bản chuyển thể không chính thức; việc
ghi tên tác giả và nhà xuất bản không có nghĩa là họ bảo trợ hoặc phê duyệt bản
dịch. Bản dịch được phát hành theo CC BY-NC-SA 4.0:
\url{https://creativecommons.org/licenses/by-nc-sa/4.0/}.
Khi sử dụng lại, phải ghi công tác giả và nguồn, dẫn giấy phép, nêu các thay đổi,
không sử dụng cho mục đích thương mại, và chia sẻ bản chuyển thể theo cùng giấy
phép hoặc giấy phép tương thích. Nội dung của bên thứ ba có điều kiện sử dụng
riêng được xử lý theo các điều kiện đó.

\noindent Thông tin xuất bản của bản gốc: sách có thể được mua với mức chiết
khấu theo số lượng cho doanh nghiệp hoặc hoạt động quảng bá bán hàng. Để biết
thêm thông tin, liên hệ \nolinkurl{special_sales@mitpress.mit.edu}.

\paragraph{Dữ liệu biên mục của Thư viện Quốc hội Hoa Kỳ.}
Spivak, David I., sinh năm 1978, tác giả.
\emph{Category theory for the sciences / David I. Spivak.}
Dòng mô tả hình thức trong dữ liệu biên mục gốc: ``trang; cm''.
Sách có tài liệu tham khảo và chỉ mục.
ISBN 978-0-262-02813-4 (bìa cứng, giấy kiềm).
Chủ đề: 1. Khoa học - Mô hình toán học. 2. Phạm trù (Toán học).
I. Nhan đề.
\texttt{Q175.32.M38S65 2014}; \texttt{512'.62-dc23}; \texttt{2014007215}.
\noindent Dòng kiểm soát lần in trong bản gốc: \texttt{10 9 8 7 6 5 4 3}.
\par\endgroup

\chapter*{Lời cảm ơn}
\addcontentsline{toc}{chapter}{Lời cảm ơn}
% Source: printed p. vii; PDF p. 8.
Tôi muốn bày tỏ lòng biết ơn sâu sắc đối với nhiều nhà khoa học mà tôi đã làm
việc cùng trong sáu năm qua. Mọi việc bắt đầu với Paea LePendu, người đầu tiên
dạy tôi về cơ sở dữ liệu khi tôi còn ngây thơ đi gõ cửa từng phòng ở khoa khoa
học máy tính của Đại học Oregon. Cuốn sách này sẽ không bao giờ được viết nếu
Tristan Nguyen và Dave Balaban không chú ý đến công trình của tôi và khuyến
khích tôi tiếp tục. Dave Balaban và Peter Gates đã là những cộng sự khoa học
của tôi ngay từ đầu; họ cố gắng hiểu điều tôi mang đến, đồng thời cũng nỗ lực
không kém để giúp tôi hiểu tất cả những gì mình còn thiếu. Peter Gates đã làm
cho hiểu biết của tôi về dữ liệu sâu sắc hơn theo nhiều cách.

Tôi cũng vô cùng may mắn được quen Haynes Miller, người đã giúp tôi có thể đến
làm việc ổn định tại MIT, với sự hỗ trợ của Clark Barwick và Jacob Lurie. Tôi
biết MIT sẽ là nơi tốt nhất trên thế giới để theo đuổi loại nghiên cứu này, và
nơi đây luôn đáp ứng kỳ vọng đó. Tôi rất vui khi làm việc với những nhà nghiên
cứu như Markus Buehler và các học viên sau đại học của anh, Tristan Giesa và Dieter
Brommer. Những ví dụ về khoa học vật liệu xuất hiện rải rác trong cuốn sách là
bằng chứng cho thấy công việc chung của chúng tôi đã ảnh hưởng đến cách suy
nghĩ của tôi nhiều đến mức nào.

Tôi cũng muốn cảm ơn các cộng sự và những người đã cùng tôi trao đổi về các
chủ đề được trình bày trong sách. Ngoài những người đã kể trên, họ còn gồm
Steve Awodey, Allen Brown, Adam Chlipala, Carlo Curino, Dan Dugger, Henrik
Forssell, David Gepner, Jason Gross, Bob Harper, Ralph Hutchison, Robert Kent,
Jack Morava, Scott Morrison, David Platt, Joey Perricone, Dylan Rupel, Guarav
Singh, Sam Shames, Nat Stapleton, Patrick Schultz, Ka Yu Tam, Ryan Wisnesky,
Jesse Wolfson và Elizabeth Wood.

Tôi xin cảm ơn Peter Kleinhenz và Peter Gates đã đọc một phiên bản trước của
cuốn sách và đưa ra những nhận xét vô cùng quý giá trước khi tôi bắt đầu dạy
lớp 18-S996 tại MIT vào mùa xuân năm 2013. Đặc biệt, hình đầu tiên trong sách,
Hình 1.1, là một phiên bản được chỉnh sửa nhẹ của sơ đồ mà Gates gửi cho tôi
để giúp giải thích động cơ của cuốn sách cho các nhà khoa học. Tôi cũng đặc
biệt cảm ơn người chấm bài của khóa 18-S996, Darij Grinberg. Anh không chỉ là
người chấm bài tốt nhất mà tôi từng có trong 14 năm giảng dạy, mà còn góp ý
cho chính cuốn sách nhiều hơn bất kỳ ai khác.

% Source: printed p. viii; PDF p. 9.
Tôi muốn cảm ơn những sinh viên trong lớp 18-S996 tại MIT đã tìm ra lỗi đánh
máy, chỉ ra các cách giải thích chưa rõ và nói chung đã giúp cải thiện cuốn
sách theo nhiều cách: Aaron Brookner, Leon Dimas, Dylan Erb, Deokhwan Kim,
Taesoo Kim, Owen Lewis, Yair Shenfeld và Adam Strandberg, cùng những người
khác. Những người ngoài lớp, V. Galchin, K. Hofmeyr, D. McAdams, D. Holmes,
C. McNally, P. O'Neill và R. Harper, cũng góp phần tìm lỗi và cải thiện sách.

Tôi cũng xin cảm ơn Marie Lufkin Lee, Marc Lowenthal, Katherine Almeida và
mọi người khác tại MIT Press đã giúp chuẩn bị cuốn sách để xuất bản. Cảm ơn
Laura Baldwin, người đã giúp tôi xử lý một số vấn đề LaTeX rất phiền phức.
Cuốn sách chắc chắn tốt hơn nhiều so với bản tôi gửi ban đầu. Tôi cũng trân
trọng sự sẵn lòng của nhà xuất bản trong việc phối hợp với tôi để đưa một
bản sao của cuốn sách ra công chúng.

Tôi cũng cảm ơn người thầy Peter Ralston, người đã dạy tôi không ngừng đặt câu
hỏi về những điều hiển nhiên. Khả năng cam kết thực hiện một dự án như thế này
và theo đuổi nó đến khi có kết quả chắc chắn đã được nâng cao sau thời gian
tôi học với ông.

Cuối cùng, tôi xin ghi nhận và bày tỏ lòng biết ơn đối với sự hỗ trợ của Văn
phòng Nghiên cứu Hải quân và Văn phòng Nghiên cứu Khoa học Không quân,
\footnote{Các mã tài trợ: N000140910466, N000141010841, N000141310260,
FA9550-14-1-0031.}
mà nếu không có sự hỗ trợ đó thì cuốn sách này hoàn toàn không thể ra đời.
Tôi tin rằng tài trợ cho nghiên cứu cơ bản là một cách tuyệt vời để bảo đảm
Hoa Kỳ tiếp tục giữ vai trò dẫn đầu thế giới trong những năm tới.
